\section{Preliminaries}

\subsection{Theta functions}

We recall some classical results for theta functions.
We will write
\begin{gather*}
\ta(x;p)=\prod\limits_{j=0}^\infty\big(1-p^jx\big)\big(1-p^{j+1}/x\big),
\end{gather*}
where $p$ is a complex number with $|p|<1$.
We refer to the case
 $p=0$, $\ta(x;0)=1-x$, as the \emph{trigonometric case}.
We will also use the shorthand notation
\begin{gather*}
\theta(x_1,\dots,x_n;p)=\ta(x_1;p)\dotsm\ta(x_n;p),
\end{gather*}
\begin{gather*}
\theta\big(x_1y^{\pm},\dots,x_ny^{\pm};p\big)=\ta(x_1y;p)\ta(x_1/y;p)\dotsm\ta(x_ny;p)
\ta(x_n/y;p).
\end{gather*}
In Section~\ref{gmis}, we will suppress the dependence on $p$ and simply write
$\theta(x)=\theta(x;p)$.


Among the properties of $\theta$ we mention the inversion formula
\begin{gather}
\label{ti}\theta(1/x;p)=-\frac1 x \theta(x;p),
\end{gather}
the quasi-periodicity
\begin{gather}
\label{qpt}\ta(px;p)=-\frac1 x \theta(x;p)
\end{gather}
and Weierstrass' addition formula
\begin{gather}\label{tadd}
\ta(xy,x/y,uv,u/v;p)-\ta(xv,x/v,uy,u/y;p)=\frac uy \ta(yv,y/v,xu,x/u;p).
\end{gather}
We will also need the identity
\begin{gather}\label{cpf}
\sum_{l=1}^k\frac{a_l\prod\limits_{j=1}^{k-2}\theta\big(a_lb_j^\pm;p\big)}
{\prod\limits_{j=1,\,j\neq l}^k\theta\big(a_la_j^\pm;p\big)}=0, \qquad
k\geq 2.
\end{gather}
To our knowledge, it was first obtained by Gustafson
\cite[Lemma~4.14]{gu}, see~\cite{r} for further references and comments.
